\documentclass[a4paper,12pt]{article}
\usepackage[utf8]{inputenc}
\usepackage[spanish]{babel}
\usepackage{amsmath}
\usepackage{geometry}

\geometry{margin=2.5cm}

\title{Solucionario: Examen de Física I - Forma A}
\date{}

\begin{document}
\maketitle

\section{Parte I: Teoría}
\begin{enumerate}
    \item \textbf{Interdisciplinariedad}: Matemáticas, Astronomía, Biología.
    \item \textbf{Cifras Significativas}: Ceros a la izquierda no son significativos. Ceros intermedios siempre son significativos.
    \item \textbf{Divisiones}: Física Clásica y Moderna. Ramas: Mecánica, Termodinámica, Óptica.
    \item \textbf{Unidades SI}: Longitud (m), Masa (kg), Tiempo (s), Temperatura (K), Intensidad (A).
\end{enumerate}

\section{Parte II: Práctica}
\begin{enumerate}
    \item \textbf{9.5 km a pies}: $9.5 \times 1000 \times 3.28084 = 31,167.98$ ft.
    \item \textbf{12 mi/h a m/s}: $12 \times 1609.34 / 3600 = 5.36$ m/s.
    \item \textbf{Notación Científica}: $(4.5 \times 1.8) \times 10^{17} = 8.1 \times 10^{17}$.
    \item \textbf{Cifras Significativas}: $15.482 + 3.1 + 102.55 = 121.132 \rightarrow 121.1$.
\end{enumerate}

\end{document}
